\section*{Summary}
\addcontentsline{toc}{section}{Summary}

\noindent\fbox{\parbox{\dimexpr\textwidth-2\fboxsep-2\fboxrule}{%
\textbf{Proposed Weisshorn nose cone}\\[3pt]
\begin{tabular}{@{}p{2.6cm}>{\raggedright\arraybackslash}p{12.8cm}@{}}
Profile & Von K\'arm\'an (Haack $C=0$), tangent to the body \\
Length & $L = \FINL$ ($L/D = 4.5$) to the theoretical apex; \FINLPHYS{} physical \\
Diameter & \SI{180}{mm} \\
Tip & spherical, $r_n = \SI{4.5}{mm}$ ($r_n/R = 0.05$) \\
Fallback & same profile at $L/D=4$ (\SI{720}{mm}) if the \SI{62}{mm} cannot be found in the \SI{5.4}{m} budget \\
\end{tabular}}}

\vspace{6pt}
\noindent The Von K\'arm\'an profile is the member of the Haack family of noses that, according to
slender-body theory, gives the lowest supersonic pressure drag for a given length and base
diameter (Section~\ref{sec:profiles}). $L/D$ is the \emph{fineness ratio}, the nose length divided
by the body diameter. Symbols and abbreviations are listed in the Nomenclature. The report covers
the FD scope (profile, length, tip radius, drag, apogee, stability); structural loads, heating,
materials and the mass budget are left to ST.

\vspace{4pt}
\noindent Why:
\begin{itemize}[leftmargin=*,itemsep=2pt]
\item \textbf{Drag.} Above Mach 1 a shock wave forms ahead of the nose, and the overpressure it
creates on the nose surface produces a pressure drag called \emph{wave drag}. Wave drag is the
only part of the vehicle drag that depends on the nose shape (Section~\ref{sec:physics}). Seven
profiles (cone, tangent ogive, parabola, power $\tfrac12$ and
$\tfrac34$, LV-Haack, Von K\'arm\'an; Section~\ref{sec:catalogue}) were compared with an
axisymmetric Euler solver, i.e.\ a CFD code for inviscid compressible flow (validated to $<0.5\,\%$
against Taylor--Maccoll, the exact solution for a cone, and grid-converged to
0.5\,\%). The solver was coupled to a trajectory model calibrated on the Launch Vehicle Overview (LVO). At $L/D = 4.5$, Von K\'arm\'an has
the lowest wave drag at every computed Mach number. It flies \SIrange{18}{22}{m} higher than the
power noses, \SIrange{69}{79}{m} higher than ogive, parabola and LV-Haack, and \SI{132}{m}
higher than the cone. It stays first in all seven robustness scenarios. At $-\SI{143}{m/kg}$ (apogee
lost per kilogram of dry mass), the margin over the ogive is worth about \SI{0.55}{kg} of dry mass.
\item \textbf{Length.} $L/D = 4.5$ adds \TRADELEN{} over the LVS length (the \SI{720}{mm} nose of the
Launch Vehicle Sizing) for \SI{62}{mm} of nose
bay, with diminishing returns beyond that: wave drag falls as $1/(L/D)^2$, while the friction and
the mass of the nose grow linearly with its length.
\item \textbf{Tip.} A \SI{4.5}{mm} spherical tip replaces the sharp apex, which cannot be made, for
\BLUNTCOST{} of wave drag (\BLUNTAPOA{} of apogee).
\item \textbf{Integration and stability.} The shoulder is tangent (no change of slope where the nose
meets the body tube); \FINTHREEU{} of the 3U payload fits inside the curved part. The static margin, the distance from the centre of
gravity to the centre of pressure in body diameters (calibres, cal), changes by \FINDSM{} cal,
which is negligible.
\item \textbf{Deliverables.} Profile coordinates, watertight STL solids of all seven designs and
of the final nose, a STEP solid of the final nose, the full Euler data set, six numbered
recommendations with owners
(Section~\ref{sec:recommendations}) and the alternatives to use if the priorities change.
\end{itemize}

\noindent\textbf{Not found in the ERT documentation} (details in Section~\ref{sec:flags}): any
Weisshorn nosecone requirement, the Weisshorn fin geometry and mass breakdown, and any supersonic
flight data.
